\exercisesection{}

\begin{enumerate}
\def\labelenumi{\arabic{enumi})}
\tightlist
\item
  Let \(E = \ell_{\mathbf{C}}^{2}(\mathbf{N})\) and let \(u \in \mathcal{L}(\mathrm{E})\) be the endomorphism such that \(u(x_{0}, \ldots, x_{n}, \ldots) = (0, x_{0}, \ldots, x_{n}, \ldots)\) for all \((x_{n})_{n \in \mathbf{N}}\) in E. Let \((e_{n})\) be the canonical orthonormal basis of E.
\end{enumerate}

\begin{enumerate}
\def\labelenumi{\alph{enumi})}
\item
  For every \(v \in \mathcal{L}(\mathrm{E})\) such that \(\|v\| < 1/2\) , the vector \(e_{0}\) is not adherent to the image of \(u + v\) .
\item
  The open ball in \(\mathcal{L}(\mathrm{E})\) of radius 1/2 centered at u contains no diagonalizable endomorphism of E.
\end{enumerate}

\begin{enumerate}
\def\labelenumi{\arabic{enumi})}
\setcounter{enumi}{1}
\item
  Let E be a complex Hilbert space. Prove directly that an endomorphism \(u\) of E is a Fredholm endomorphism if and only if \(\pi(u)\) is invertible in \(\mathcal{C}alk(E)\) (Atkinson's theorem). (Use the spectral properties of compact operators, cf. § 1 of IV, p. 146.)
\item
  Let E be a complex Hilbert space and let \(\mathrm{B}=(e_{i})_{i\in\mathrm{I}}\) be an orthonormal basis of E. For every \(i\in I\) , we denote by \(p_{i}\) the orthogonal projector with image \(Ce_{i}\) .
\end{enumerate}

\begin{enumerate}
\def\labelenumi{\alph{enumi})}
\item
  Let K be a precompact subset of E. For every \(\varepsilon > 0\) there exists a finite subset J of I such that, for every \(x \in \mathbf{K}\) , \(\|x - \sum_{i \in \mathbf{J}} p_j(x)\| \leqslant \varepsilon\) .
\item
  Let \(u \in \mathcal{D}_{\mathrm{B}}(\mathrm{E})\) and let us denote by \(\lambda = (\lambda_{i})_{i \in \mathrm{I}}\) its family of eigenvalues. The family \((\lambda_{i} p_{i})\) is summable in \(\mathcal{L}(\mathrm{E})\) endowed with the topology of precompact convergence and whose sum is equal to u (cf. prop. 1 of IV, p. 147)
\end{enumerate}

\begin{enumerate}
\def\labelenumi{\arabic{enumi})}
\setcounter{enumi}{3}
\item
  Let E be a complex Hilbert space of infinite dimension. There exists a compact endomorphism of E that is not of finite trace.
\item
  \begin{enumerate}
  \def\labelenumii{\alph{enumii})}
  \tightlist
  \item
  Let \(E = L^{2}([0,1], \mu)\) where \(\mu\) is the Lebesgue measure. The image A of the algebra \(L^{\infty}([0,1], \mu)\) by the map that associates to f the endomorphism of multiplication by f is a maximal commutative subalgebra of \(\mathcal{L}(E)\) which is not of the form \(\mathcal{D}_{\mathrm{B}}(\mathrm{E})\) for an orthonormal basis B of E. (Verify that the algebra A contains no minimal projector.)
  \end{enumerate}
\end{enumerate}

\begin{enumerate}
\def\labelenumi{\alph{enumi})}
\setcounter{enumi}{1}
\tightlist
\item
  Let E be an infinite-dimensional Hilbert space. There exists a maximal commutative subalgebra of \(\mathcal{L}(\mathrm{E})\) which is not of the form \(\mathcal{D}_{\mathrm{B}}(\mathrm{E})\) for B an orthonormal basis of E.
\end{enumerate}

\begin{enumerate}
\def\labelenumi{\arabic{enumi})}
\setcounter{enumi}{5}
\tightlist
\item
  Let E be a complex Hilbert space.
\end{enumerate}

\begin{enumerate}
\def\labelenumi{\alph{enumi})}
\item
  Let u be a compact normal endomorphism of E. Every closed subspace of E invariant under u is invariant under \(u^{*}\) .
\item
  If E is infinite-dimensional, there exist a normal endomorphism u of E and a closed subspace \(F \subset E\) such that:
\end{enumerate}

\begin{enumerate}
\def\labelenumi{(\roman{enumi})}
\item
  One has \(u(\mathrm{F}) \subset \mathrm{F}\) ;
\item
  The space F is not invariant under \(u^{*}\) ;
\item
  The space F° is not invariant under u ;
\item
  The endomorphism of F induced by \(u\) by passing to subspaces is not normal.
\end{enumerate}

\begin{enumerate}
\def\labelenumi{\arabic{enumi})}
\setcounter{enumi}{6}
\tightlist
\item
  The goal of this exercise is to give a more direct proof of Theorem 1 of IV, p. 149. Let E be an infinite-dimensional complex Hilbert space and let u be a compact and normal endomorphism of E.
\end{enumerate}

\begin{enumerate}
\def\labelenumi{\alph{enumi})}
\item
  Show that, for every nonzero complex number \(\lambda\), the eigenspace \(\operatorname{Ker}(u - \lambda 1_{\mathrm{E}})\) is of finite dimension.
\item
  Suppose that \(u\) is Hermitian. Show that there exists \(x_0 \neq 0\) in E such that
\end{enumerate}

\[
\frac {\| u (x _ {0}) \| ^ {2}}{\| x _ {0} \| ^ {2}} = \sup _ {x \in \mathrm{E} - \{0 \}} \frac {\| u (x) \| ^ {2}}{\| x \| ^ {2}}.
\]

\begin{enumerate}
\def\labelenumi{\alph{enumi})}
\setcounter{enumi}{2}
\item
  Show that \(x_{0}\) is an eigenvector of \(u\) for an eigenvalue \(\lambda\) such that \(|\lambda| = \|u\|\) .
\item
  Show that there exists a countable set I, an orthonormal family \((e_{i})_{i\in\mathrm{I}}\) in E and a family \((\lambda_{i})_{i\in\mathrm{I}}\) of real numbers such that
\end{enumerate}

\[
u (x) = \sum_ {i \in I} \lambda_ {i} \langle e _ {i} | x \rangle e _ {i}
\]

for all \(x \in E\) .

\begin{enumerate}
\def\labelenumi{\alph{enumi})}
\setcounter{enumi}{4}
\tightlist
\item
  Extend the result to the case where u is compact and normal. (Write \(u = u_{1} + iu_{2}\) where \(u_{1}\) and \(u_{2}\) are compact, Hermitian, and mutually commuting.)
\end{enumerate}

\begin{enumerate}
\def\labelenumi{\arabic{enumi})}
\setcounter{enumi}{7}
\tightlist
\item
  We endow the group R/Z with its normalized Haar measure \(\mu\) .
\end{enumerate}

\begin{enumerate}
\def\labelenumi{\alph{enumi})}
\tightlist
\item
  Let \(h \in \mathcal{L}^{2}(\mathbf{R}/\mathbf{Z})\) . For \(f \in \mathcal{L}^{2}(\mathbf{R}/\mathbf{Z})\) we define a function \(\mathrm{C}_{h}(f)\) on R/Z by the formula
\end{enumerate}

\[
\mathrm{C} _ {h} (f) (x) = \int_ {\mathbf {R} / \mathbf {Z}} f (t) h (x - t) d \mu (t).
\]

The map \(f \mapsto \mathrm{C}_h(f)\) is a continuous endomorphism of \(\mathcal{L}^2(\mathbf{R}/\mathbf{Z})\) , and by passing to quotients defines a compact normal endomorphism of \(\mathrm{L}^2(\mathbf{R}/\mathbf{Z})\) .

\begin{enumerate}
\def\labelenumi{\alph{enumi})}
\setcounter{enumi}{1}
\item
  Determine the spectrum of \(C_{h}\) and the spectral multiplicity of the elements of this spectrum. (Expand h in a Fourier series.)
\item
  There exists \(h \in \mathcal{C}(\mathbf{R}/\mathbf{Z})\) such that \(C_{h}\) is not of finite trace.
\end{enumerate}

\begin{enumerate}
\def\labelenumi{\arabic{enumi})}
\setcounter{enumi}{8}
\tightlist
\item
  We denote by E the complex Hilbert space \(L^{2}([0,1])\) of functions that are square-integrable on \([0,1]\) with respect to Lebesgue measure \(\mu\) .
\end{enumerate}

Let \(u_{0} \in \mathcal{L}(E)\) be the Volterra endomorphism on E, such that

\[
u _ {0} (f) (x) = \int_ {[ 0, x ]} f (t) d \mu (t)
\]

for all \(f \in \mathrm{L}^2([0,1])\) and almost all \(x \in [0,1]\) .

\begin{enumerate}
\def\labelenumi{\alph{enumi})}
\item
  The endomorphism \(u_{0}\) is compact and \(\operatorname{Sp}(u_{0})=\{0\}\) .
\item
  The endomorphism \(u_0^* u_0\) has finite trace; compute its trace and deduce the value of the series
\end{enumerate}

\[
\sum_ {n = 1} ^ {+ \infty} \frac {1}{n ^ {2}}.
\]

\begin{enumerate}
\def\labelenumi{\alph{enumi})}
\setcounter{enumi}{2}
\tightlist
\item
  Let \(v_{0}\) be the rank-1 endomorphism of E defined by
\end{enumerate}

\[
v _ {0} (f) = \left(\int_ {0} ^ {1} f (t) d \mu (t)\right) \cdot 1,
\]

where 1 denotes the class of the constant function 1. Let \(u = u_0 - v_0\) . Show that \(u\) is a Hilbert--Schmidt operator, and compute the kernel N such that

\[
u (f) (x) = \int_ {0} ^ {1} \mathrm{N} (x, y) f (y) d \mu (y)
\]

for \(f \in \mathrm{E}\) and almost all \(x \in [0,1]\) .

\begin{enumerate}
\def\labelenumi{\alph{enumi})}
\setcounter{enumi}{3}
\item
  Determine the eigenvalues of u and their spectral multiplicity. (Compute \(u(e_{n})\) , where \(e_{n}(t) = e^{2i\pi nt}\) for all \(n \in Z\) .)
\item
  Compute the trace of \(u^{*}u\) and deduce anew the value of the series whose general term is \(1/n^{2}\) for \(n \geqslant 1\) .
\item
  Generalize these arguments to show that, if \(k \geqslant 1\) is a strictly positive integer, then
\end{enumerate}

\[
\sum_ {n = 1} ^ {+ \infty} \frac {1}{n ^ {2 k}} \in \pi^ {2 k} {\bf Q} ^ {*}.
\]

\begin{enumerate}
\def\labelenumi{\arabic{enumi})}
\setcounter{enumi}{9}
\tightlist
\item
  Let E be a complex Hilbert space. Let A be a star-subalgebra of \(\mathcal{L}(\mathrm{E})\) . Suppose that the representation of A in E is irreducible and that A contains a nonzero compact endomorphism.
\end{enumerate}

\begin{enumerate}
\def\labelenumi{\alph{enumi})}
\item
  The star-subalgebra A contains a nonzero finite-rank projector. (Show that A contains a nonzero compact Hermitian endomorphism and apply the functional calculus.)
\item
  The star-subalgebra A contains a rank-1 projector. (If \(p \in \mathrm{A}\) is a nonzero projector of minimal rank, show that the algebra \(p\mathrm{A}p\) is contained in \(\mathbf{C}p\) .)
\item
  It follows that A contains \(\mathcal{L}^{\mathrm{c}}(\mathrm{E})\) .
\item
  Conversely, every star-subalgebra of \(\mathcal{L}(\mathrm{E})\) containing \(\mathcal{L}^{\mathrm{c}}(\mathrm{E})\) is irreducible.
\end{enumerate}

\begin{enumerate}
\def\labelenumi{\arabic{enumi})}
\setcounter{enumi}{10}
\tightlist
\item
  Let E be a complex Hilbert space. Let \(A \subset \mathcal{L}(\mathrm{E})\) be a set of compact endomorphisms such that \(u \circ v = v \circ u\) for every pair \((u, v)\) of elements of A and \(u^{*} \in A\) for all \(u \in A\) .
\end{enumerate}

Show that there exists a countable orthonormal family \((e_i)_{i\in \mathrm{I}}\) and a family \((\Lambda_i)_{i\in \mathrm{I}}\) of mappings \(\mathrm{A}\to \mathbf{C}\) such that

\[
u (x) = \sum_ {i \in \mathrm{I}} \Lambda_ {i} (u) \langle e _ {i} | x \rangle e _ {i}
\]

for every element u of A and every x in E. If A is a unital subalgebra of \(\mathcal{L}(\mathrm{E})\) , then \(\Lambda_{i}\) is a character of A for every \(i \in I\) .

\begin{enumerate}
\def\labelenumi{\arabic{enumi})}
\setcounter{enumi}{11}
\tightlist
\item
  If E is a complex Hilbert space and u is a compact normal endomorphism of E, show that u has a cyclic vector (cf. def. 3 of IV, p. 191) if and only if the following conditions are satisfied:
\end{enumerate}

\begin{enumerate}
\def\labelenumi{(\roman{enumi})}
\item
  The kernel of u is reduced to 0;
\item
  The spectral multiplicity of every nonzero eigenvalue of u is equal to 1.
\end{enumerate}

\begin{enumerate}
\def\labelenumi{\arabic{enumi})}
\setcounter{enumi}{12}
\tightlist
\item
  Let E be a complex Hilbert space of countable type, let \((e_{n})_{n\in\mathbf{N}}\) be an orthonormal basis of E, and let u be an endomorphism of E. Show that u is compact if and only if one has
\end{enumerate}

\[
\lim _ {n \to + \infty} \sup _ {x \in \{e _ {0}, \dots , e _ {n - 1} \} ^ {\circ} - \{0 \}} \frac {\| u (x) \| ^ {2}}{\| x \| ^ {2}} = 0.
\]

\begin{enumerate}
\def\labelenumi{\arabic{enumi})}
\setcounter{enumi}{13}
\item
  Let \(\mathrm{X} = [0,1]\) be equipped with the Lebesgue measure, and let \(\mathrm{N}\colon \mathrm{X}\times \mathrm{X}\to \mathbf{R}_{+}\) be such that \(\mathrm{N}(x,y) = |x - y|\) . The endomorphism \(u_{\mathrm{N}}\) of \(\mathrm{L}^2 (\mathrm{X})\) is not positive.
\item
  Let X be a locally compact topological space, let \(\mu\) be a positive bounded measure on X with support X, and let \(N \in \mathcal{C}_{b}(X \times X)\) .
\end{enumerate}

\begin{enumerate}
\def\labelenumi{\alph{enumi})}
\tightlist
\item
  There exists an endomorphism \(u\) of \(\mathcal{C}_b(\mathrm{X})\) such that
\end{enumerate}

\[
u (f) (y) = \int_ {\mathrm{X}} \mathrm{N} (x, y) f (x) d \mu (x)
\]

for every function \(f \in \mathcal{C}_b(\mathrm{X})\) and every \(y \in \mathrm{X}\) . In what follows, we assume that \(u\) is compact; this is the case, for example, if the space \(\mathrm{X}\) is compact.

\begin{enumerate}
\def\labelenumi{\alph{enumi})}
\setcounter{enumi}{1}
\item
  The image of \(u_{\mathrm{N}}\) is contained in \(\mathcal{C}_b(\mathrm{X})\) and \(u_{\mathrm{N}}\) defines a continuous linear map from \(\mathrm{L}_{\mathbf{C}}^{2}(\mathrm{X})\) into \(\mathcal{C}_b(\mathrm{X})\) .
\item
  Let \(f \in \mathrm{L}_{\mathbf{C}}^{2}(\mathrm{X})\) and \(\lambda \in \mathbf{C}^*\) such that \(u_{\mathrm{N}}(f) = \lambda f\) . Then \(f \in \mathcal{C}_b(\mathrm{X})\) .
\end{enumerate}

\(d)\) There exists a countable family \((f_i)_{i\in I}\) in \(\mathcal{C}_b(\mathrm{X})\) and a family \((\lambda_i)_{i\in I}\) of nonzero real numbers such that for all \(f\in \mathrm{L}_{\mathbf{C}}^2 (\mathrm{X})\) , we have

\[
u (f) = \sum_ {i \in I} \lambda_ {i} \langle f _ {i} | f \rangle f _ {i}
\]

in \(\mathcal{C}_b(\mathrm{X})\) .

\begin{enumerate}
\def\labelenumi{\alph{enumi})}
\setcounter{enumi}{4}
\tightlist
\item
  Suppose that \(u_{\mathrm{N}}\) is a positive endomorphism of \(\mathrm{L}_{\mathbf{C}}^{2}(\mathrm{X})\) . For every \(i \in \mathrm{I}\) , we denote by \(h_i \in \mathcal{C}_b(\mathrm{X} \times \mathrm{X})\) the function defined by \(h_i(x,y) = \overline{f_i(x)} f_i(y)\) . Then
\end{enumerate}

\[
N = \sum_ {i \in I} \lambda_ {i} h _ {i}
\]

in \(\mathcal{C}(\mathrm{X}\times\mathrm{X})\) endowed with the topology of compact convergence. (Prove that the series converges using Dini's theorem (TG, X, p. 34, th. 1), then prove the equality.)

\begin{enumerate}
\def\labelenumi{\arabic{enumi})}
\setcounter{enumi}{15}
\tightlist
\item
  Let \(\mathrm{X} = [0,1]\) endowed with the Lebesgue measure.
\end{enumerate}

\begin{enumerate}
\def\labelenumi{\alph{enumi})}
\item
  For any sequence \((f_n)_{n\in \mathbf{N}}\) of continuous functions on X that is orthonormal in \(\mathrm{L}_{\mathbf{C}}^2 (\mathrm{X})\) , there exists a kernel \(k\in \mathcal{C}(\mathrm{X}\times \mathrm{X})\) such that
\item
  The endomorphism \(u_{k}\) of \(\mathrm{L}_{\mathbf{C}}^{2}(\mathrm{X})\) is positive;
\item
  For every \(n \in \mathbf{N}\) , the function \(f_n\) is an eigenfunction of \(u_k\) whose eigenvalue is nonzero;
\item
  The family \((f_n)\) is a basis of the image of the endomorphism \(u_k\) .
\item
  There exist kernels \(k \in \mathcal{C}(\mathrm{X} \times \mathrm{X})\) and orthonormal families of eigenfunctions \((f_n)\) of the endomorphism \(u_k\) of \(\mathrm{L}_{\mathbf{C}}^2(\mathrm{X})\) consisting of bounded functions, such that the sequence \((\|f_n\|_\infty)\) is not bounded.
\end{enumerate}

\begin{enumerate}
\def\labelenumi{\arabic{enumi})}
\setcounter{enumi}{16}
\item
  Let E and F be Hilbert spaces. The bilinear map \((u,v)\mapsto\) \(\mathrm{Tr}(uv)\) defines an isometric isomorphism from the Banach space \(\mathcal{L}_1(\mathrm{E};\mathrm{F})'\) into the space \(\mathcal{L}(\mathrm{F};\mathrm{E})\) .
\item
  Let E be a complex Hilbert space and let \(I_{E}\) be the set defined in no. 3 of IV, p. 151.
\end{enumerate}

Let \(n \in I_{E}\) . Let F be a complex Hilbert space. The map \(u \mapsto \lambda_{n}(|u|)\) from \(\mathcal{L}^{c}(E;F)\) into \(R_{+}\) is continuous.

\begin{enumerate}
\def\labelenumi{\arabic{enumi})}
\setcounter{enumi}{18}
\item
  We keep the notations of Lemma 5 of IV, p. 161. Show that if the image of \(f\) is contained in \(\mathbf{R}\) , if \(f\) is strictly convex, and if \(\varrho_n > 0\) , then there is equality in inequality (6) if and only if \(a_i = b_i\) for \(0 \leqslant i \leqslant n\) .
\item
  Let E be a pre-Hilbert space, let F be a Banach space, and let \(u: E \to F\) be a linear map.
\end{enumerate}

\begin{enumerate}
\def\labelenumi{\alph{enumi})}
\item
  Prove that u is a compact linear map if and only if \(\|u(e_{n})\|\to0\) for every orthonormal sequence \((e_{n})\) of E. (To prove that this condition is sufficient, use TG, IX, p. 20, prop. 14.)
\item
  Suppose that E is a Hilbert space of countable type and that F = E. For there to exist an orthonormal basis \((e_{n})_{n\in\mathbf{N}}\) of E such that \(\|u(e_{n})\|\to0\) , it is necessary and sufficient that u not be a Fredholm endomorphism.
\item
  Suppose that E is a Hilbert space and that F = E. If one has \(\langle u(e_{n}) | e_{n} \rangle \to 0\) for every orthonormal sequence \((e_{n})\) of E, then u is a compact linear map.
\end{enumerate}

\begin{enumerate}
\def\labelenumi{\arabic{enumi})}
\setcounter{enumi}{20}
\tightlist
\item
  Let E be a Hilbert space of countable type.
\end{enumerate}

\begin{enumerate}
\def\labelenumi{\alph{enumi})}
\item
  Let A be an infinite set of unit vectors of E which converges weakly to 0 along the filter of complements of finite sets. Prove that there exists a sequence \((x_{n})_{n\geqslant1}\) with values in A and an orthonormal sequence \((e_{n})_{n\geqslant1}\) in E such that \(\|x_{n}-e_{n}\|\leqslant1/n\) for all \(n\geqslant1\) .
\item
  Let u be a linear map from E into E, not necessarily continuous. Suppose that for every orthonormal sequence \((e_{n})\) in E, the sequence \((u(e_{n}))\) converges weakly to 0. Prove that u is continuous and compact.
\item
  Let u be an endomorphism of E. There exists an orthonormal basis \((e_{n})_{n\in\mathbf{N}}\) of E such that the sequence \((u(e_{n}))_{n\in\mathbf{N}}\) converges weakly to 0 if and only if there does not exist an endomorphism v of E such that \(v\circ u-1_{E}\) is compact.
\end{enumerate}

\begin{enumerate}
\def\labelenumi{\arabic{enumi})}
\setcounter{enumi}{21}
\tightlist
\item
  Let E be a complex Hilbert space. Let u and v be positive endomorphisms of E.
\end{enumerate}

\begin{enumerate}
\def\labelenumi{\alph{enumi})}
\tightlist
\item
  For every real number p such that \(0 < p < 1\) and every \(x \in R_{+}\) , we have
\end{enumerate}

\[
x ^ {p} = \frac {\sin (p \pi)}{\pi} \int_ {0} ^ {+ \infty} t ^ {p} \left(\frac {1}{t} - \frac {1}{t + x}\right) d t.
\]

\begin{enumerate}
\def\labelenumi{\alph{enumi})}
\setcounter{enumi}{1}
\item
  Suppose that \(u \leqslant v\) . For every real number \(p\) such that \(0 < p < 1\) , we have \(u^p \leqslant v^p\) . (Use Lemma 14 of I, p. 119, b) and the preceding question.)\\
\item
  Let \(n \in \mathbf{N}\) . We set \(a = v^{n/2} \circ u^n \circ v^{n/2}\) and \(b = (v^{1/2} \circ u \circ v^{1/2})^n\) . If \(a \leqslant 1_{\mathrm{E}}\) , then \(b \leqslant 1_{\mathrm{E}}\) .
\item
  Suppose that u and v have finite trace. For every integer \(n \in N\) , we have
\end{enumerate}

\[
\operatorname{Tr} \left(\left(v ^ {1 / 2} \circ u \circ v ^ {1 / 2}\right) ^ {n}\right) \leqslant \operatorname{Tr} \left(v ^ {n / 2} \circ u ^ {n} \circ v ^ {n / 2}\right).
\]

(Let \(a\) and \(b\) be as in question \(c\) ); by using Lemma 5 of IV, p. 161, reduce to proving that the largest eigenvalue of \(\widehat{\Lambda}^{n}b\) is less than that of \(\widehat{\Lambda}^{n}a\) ; replace E by \(\widehat{\Lambda}^{n}\mathrm{E}\) and \(u\) , \(v\) by \(\widehat{\Lambda}^{n}u\) and \(\widehat{\Lambda}^{n}v\) , respectively, then apply the preceding question.)
% semantic-fragments: legacy-input-seam
\relax{}
